\begingroup\fontsize{10}{11.8}\selectfont\setlength{\parskip}{2pt}
\section{Conclusiones}
\label{v:conclusiones}
Las distribuciones de Bose--Einstein y Fermi--Dirac se obtienen como leyes
de equilibrio de las completaciones simétrica y exterior de los estados
con trayectoria y memoria. La exclusión algebraica de Pauli y el límite
diluido de Maxwell--Boltzmann forman parte de esa misma derivación.
La realización electromagnética de la ocupación bosónica produce,
a continuación, el espectro de Planck, el flujo de Stefan--Boltzmann y
las leyes de desplazamiento de Wien. Queda establecida así la conexión
entre estructura del estado, ocupaciones admisibles y observables
radiativos en las representaciones especificadas.

El mecanismo estadístico está determinado por la paridad composicional.
El producto exterior hace nilpotentes los operadores de creación e
idempotentes los operadores de ocupación, con espectro contenido en
\(\{0,1\}\); la completación simétrica admite ocupaciones enteras no
negativas. La segunda cuantización transporta los entrelazadores
contractivos, con el dominio adicional señalado para los morfismos
fuera de esa clase. La identificación relativista de paridad y espín
mantiene las condiciones de representación de rotaciones y localidad:
la exclusión algebraica y esa identificación tienen hipótesis distintas.

La factorización gran canónica convierte esas ocupaciones en los
denominadores característicos de Bose--Einstein y Fermi--Dirac.
La condición bosónica sobre el potencial químico garantiza la
convergencia de las trazas modales. En el régimen termodinámico, la
variación de la entropía combinatoria bajo restricciones de energía y población
proporciona una caracterización complementaria de las mismas leyes.
En la pantalla finita declarada,
los conteos microcanónicos reproducen por convolución las cardinalidades
fermiónica y bosónica mediante aritmética entera. Las cotas del régimen
diluido cuantifican la aproximación común de Maxwell--Boltzmann.

La ley de Planck resulta de componer tres datos de la realización
electromagnética libre: ocupación bosónica con potencial químico nulo,
densidad de dos modos transversales y energía por excitación.
El modo constante permanece fijado como fondo. Las estimaciones
uniformes en el origen y en la cola espectral justifican el paso de las
sumas discretas de la cavidad a las densidades continuas. Los momentos
de Bose determinan las densidades de número, energía y entropía.
La integración angular del flujo isotrópico proporciona la ley de
Stefan--Boltzmann; la extremización espectral proporciona los máximos
de Wien en frecuencia y longitud de onda.

La constante de Stefan--Boltzmann es el coeficiente de la composición
entre ocupación bosónica, densidad de modos y proyección angular del
flujo. Su expresión en unidades de calendario relaciona acción,
duración e información trítica. Las constantes de Wien caracterizan
los puntos estacionarios de dos densidades enlazadas por un cambio de
medida, cuyo jacobiano explica la diferencia entre sus máximos.
Estas construcciones distinguen los coeficientes estructurales, las
unidades dimensionales y las magnitudes experimentales con las que
se comparan.

La constante de Boltzmann actúa como transductor entre información
adimensional y entropía dimensional. La energía de decisión procede
de la acción y del período; la sección térmica especificada
individualiza el transductor y la temperatura, con sus leyes de cambio
de base. La relación térmica del calendario conserva la unidad de
acción, la unidad temporal y la pareja Boltzmann--temperatura.
El origen de la energía, la medida de las continuaciones y su lectura
térmica quedan de este modo enlazados por operaciones de tipos
diferenciados.

La cadena renovación--cociente--memoria--lazo realiza ese enlace. La renovación
$(20,24,25,23)$ y los cocientes tríticos producen el portador graduado;
la memoria positiva conserva $N-L$; el lazo $8{:}1$ fuerza
$\nu_W=11/9$ y $r=1/10$; y el agrupamiento $j=3n+b$ produce
$q=1/1000$ sin perder el residuo. Con los umbrales $23,26,29$ y las
multiplicidades $1,380,649$ se obtiene exactamente
$1{,}380649\times10^{-23}$. Las corrientes ascendentes y descendentes
fuerzan el desfase sectorial nulo, mientras el marcador restituye la
partición después de normalizar. La equivalencia centrada conserva la
respuesta térmica y el origen conserva el coste absoluto en área, masa y
radio. Finalmente,
$G_a k_B\Theta_{{\rm clk},a}\log3=c^4L_\alpha$ y la identidad del
transductor componen la salida espectral con la sección dimensional.

La dinámica que antecede a la ocupación dispone de una realización
unitaria explícita. El Hamiltoniano del calendario tiene una exponencial
que reproduce exactamente el desplazamiento; su acción variacional da
la ecuación de Schrödinger correspondiente. La formulación general
mantiene la continuidad fuerte y el dominio del generador autoadjunto.
En la realización periódica finita, el propagador de Floquet reproduce
el desplazamiento de \(108\) estados. Para los estados de base
\(\rho_{j_0}=|j_0\rangle\langle j_0|\), el vector de expectativas de
los dos observables de orden tiene periodo primitivo de \(108\) ciclos
de la excitación periódica.
La conjugación unitaria transporta
conjuntamente estado, observable y propagador. Para las perturbaciones
generales se conservan la estimación espectral y la cota temporal de
error, distinguiendo periodicidad, covariancia y estabilidad cuantificada.

\Needspace{11\baselineskip}
La interpretación de la entropía depende del objeto sobre el que se
evalúa. Para el estado fermiónico cuasilibre, conservador del número
e invariante bajo transformaciones de calibre en dimensión finita,
el correlador de un cuerpo \(0<C<I\) y su generador modular \(K_1\)
se determinan recíprocamente mediante la ley logística.
Las entropías de ocupación, de un operador de densidad y de una
reducción entrelazada están relacionadas por mapas explícitos y
conservan sus dominios respectivos. La existencia del límite
termodinámico exige la independencia respecto de la sucesión cofinal;
la condensación se caracteriza por la saturación del sector excitado.

La medida estacionaria del envolvente areal--volumétrico queda fijada
por la incidencia y la normalización. Su tasa máxima \(\log\varphi\)
se acompaña de una identidad de divergencia relativa que cuantifica el
déficit de cualquier otra medida estacionaria. La razón de pesos
\(\varphi^{-2}\) pertenece a esa ponderación de historias; la órbita
cronológica conserva su propio reparto temporal. La distinción permite
identificar la dinámica que produce cada distribución y el significado
de la entropía correspondiente.

El área total y el operador modular proporcionan dos observables
complementarios: la primera conserva la ocupación conjunta y el segundo
pondera de modo distinto las dos incidencias cuando \(\Delta_{\rm mod}\ne0\).
Con esa condición y normalización conocida, su lectura simultánea
recupera las ocupaciones sectoriales; para \(\Delta_{\rm mod}=0\),
el Hamiltoniano modular es escalar y esa recuperación no es posible.
La identidad finita expresa la variación de entropía como la suma
ponderada de variaciones areales menos un déficit de entropía relativa
no negativo. La inversión de ocupaciones
conserva la entropía y transforma el área en su complementaria; ambas
propiedades quedan determinadas sin identificar igualdad entrópica
con igualdad de estructura.

La conservación de memoria sostiene estas relaciones. El núcleo
aritmético mantiene las dos evaluaciones de suma y producto, sus
cocientes y la historia de las transiciones, distinguiendo el estado
transportado de sus observables. El acoplamiento bilateral conserva
la norma y recupera la entrada desde los dos canales; la invariancia
de otros observables obedece a la condición de conmutación establecida.
La recuperación del estado y la preservación de una lectura particular
son, por tanto, propiedades determinadas por sus operadores respectivos.

La realización de Fourier del registro posicional completa esta tesis de
conservación con dos fronteras exactas. Todo estado no nulo sobre
\(\Lambda=(\mathbb Z/27\mathbb Z)^2\) satisface un producto de soportes
no menor que \(729\), y todo estado unitario satisface
\(H_{\rm pos}+H_F\ge\log729\). Los indicadores normalizados de cuatro
subgrupos rectangulares saturan simultáneamente ambas cotas. La prueba
entrópica se obtiene por interpolación de Riesz--Thorin y diferenciación
en \(p=2\). Estas entropías de distribución conservan su dominio propio,
distinto de los dominios de entrelazamiento, von Neumann y termodinámica.
Los lectores dodecafásicos complementarios prueban, además, que dos
palabras con idéntico histograma pueden conservar correlaciones ordenadas
distintas.

El observador primo de vacancias retiene otra información que el recuento
marginal descarta: combina \(z_p\) con el peso \(\log p\) y produce
\(R_{\rm vac}(X)\). Su resolución modal es una identidad con convergencia
simétrica o de Fejér; los pares conjugados garantizan su realidad. Las seis
escalas conservadas y los treinta modos de \(X=10^8\) son testigos finitos,
con reproducción independiente hasta \(10^6\). La promoción asintótica
requiere cotas modales uniformes; la criba convencional interviene después
como lector de verificación.

La densidad fotónica queda además enlazada, sin cambiar su generador, con
el nivel regularizado \(A_2^{\mathrm{BG}}\): la ecuación funcional da
\(\mathcal Z_3(3)=4\pi_{\mathrm{HMT}}^2\log A_2^{\mathrm{BG}}\), y su
sustitución en los momentos térmicos produce exactamente la densidad
adimensional de número fotónico y las razones medias de energía y entropía.
Esta composición es la recíproca fotónica del desarrollo espectral
Bendersky--Glaisher; conserva separadas la genealogía HMT de \(\pi\), la
regularización espectral y la realización electromagnética.

La distribución completa del área y la entropía admiten preparaciones con
energías distintas. La sección~\ref{ins29:v:orientacion} localiza esa
diferencia en la información intrafibra: el intercambio sectorial conserva
el área, y su inverso alcanza el máximo de energía recuperable para el
Hamiltoniano especificado. El restituidor condicionado caracteriza además
el mínimo de energía libre a marginal fija.


El resultado conjunto es una cadena explícita de construcción:
la trayectoria y su paridad determinan las completaciones de ocupación;
el equilibrio determina Bose--Einstein y Fermi--Dirac; y la ocupación
bosónica, realizada sobre los modos electromagnéticos, determina
Planck, Stefan--Boltzmann y Wien. La acción, el calendario y la
transducción térmica enlazan sus escalas, mientras los mapas de memoria
y reducción precisan la información conservada. Las leyes estadísticas
y radiativas quedan relacionadas por esa estructura, con las
condiciones de cada composición expresadas en sus propios dominios.

\par\endgroup
